Discrete choice software has a long tradition of emphasizing model specification, estimation, and econometric reporting.
Biogeme is a Python-based reference platform for discrete choice models and is widely used for transportation and choice-modeling research \citep{bierlaire2023biogeme}.
Apollo is a flexible R package that supports a broad set of discrete choice models, including advanced structures such as mixed logit, latent class, and hybrid choice models \citep{hess2019apollo}.
TorchDCM does not replace these tools; instead, it uses them as reference implementations in a benchmark system and focuses on PyTorch-native computation.

R packages provide another important software line for random utility modeling.
The \texttt{mlogit} package provides an established interface for multinomial logit and related random utility models in R \citep{croissant2020mlogit}.
The \texttt{gmnl} package extends this ecosystem toward generalized multinomial logit and heterogeneous preference models \citep{sarrias2017gmnl}.
These packages are useful baselines because they are widely cited, stable, and attached to public example datasets.
TorchDCM therefore includes validation wrappers that compare parameters, covariance matrices, standard errors, and t-statistics against these packages on their canonical data examples.

Python choice-modeling packages increasingly target performance and machine-learning integration.
The \texttt{xlogit} package provides GPU-accelerated estimation for mixed logit models and is an important Python reference for simulation-based choice modeling \citep{arteaga2022xlogit}.
Other packages such as PyLogit and ML-oriented choice libraries broaden the Python ecosystem, but they do not provide the same combination of PyTorch-native tensors, broad econometric inference outputs, and cross-package benchmark reports targeted by TorchDCM.
The closest conceptual relationship is therefore not a single package, but the combination of econometric model interfaces, tensor differentiation, and reproducible software benchmarking.

Software papers in operations research often contribute by turning complex algorithms into open, modular, and benchmarked tools.
Recent IJOC-style software papers follow this pattern by identifying an implementation gap, documenting package design choices, and reporting reproducible computational comparisons.
TorchDCM follows the same software contribution logic: the paper is not only about a new estimator, but about an extensible implementation and a benchmark system that allows researchers to verify estimator behavior across packages.
